% ============================================================================
%  SECCIÓN 3.13: CONFIGURACIÓN DE GIT Y GITHUB
% ============================================================================

\section{Configuraci\'on de Git y GitHub}

El control de versiones es una pr\'actica fundamental en el desarrollo de
software, pues permite registrar, organizar y rastrear los cambios realizados
sobre el c\'odigo fuente \citep{gitdocs}. Para el presente proyecto se
utilizaron Git como sistema de control de versiones y GitHub como plataforma de
alojamiento remoto del repositorio \citep{githubdocs}.

El proyecto se encuentra alojado en un repositorio p\'ublico de GitHub, en la
direcci\'on:

\begin{center}
  \url{https://github.com/Joaquinfr87/proyecto-movil}
\end{center}

\subsection{Estrategia y estructura de ramas}

Para mantener un flujo de trabajo ordenado y desacoplar el c\'odigo de la
aplicaci\'on respecto a los documentos acad\'emicos, el repositorio est\'a
organizado en tres ramas principales:

\begin{table}[H]
  \centering
  \caption{Estructura de ramas del repositorio en GitHub}
  \label{tab:ramas-repositorio}
  \contenidoConFuente{%
    \begin{tablaAPA}
      \begin{tabular}{@{}p{3.2cm}p{12.4cm}@{}}
        \toprule
        \textbf{Rama} & \textbf{Prop\'osito y descripci\'on} \\
        \midrule
        \texttt{main} & Rama principal y estable del proyecto, reservada para versiones consolidadas y liberaciones. \\
        \addlinespace
        \texttt{docs} & Documentaci\'on formal del proyecto, incluyendo el informe acad\'emico en \LaTeX{} y los recursos documentales asociados. \\
        \addlinespace
        \texttt{develop} & C\'odigo fuente de la aplicaci\'on m\'ovil (sin documentaci\'on). Act\'ua como rama de integraci\'on continua donde se fusionan (\emph{merge}) las distintas ramas de caracter\'isticas y trabajo del equipo. \\
        \bottomrule
      \end{tabular}
    \end{tablaAPA}%
  }{Elaboraci\'on propia en base a la estrategia de versionado del equipo.}
\end{table}

\subsection{Configuraci\'on inicial del repositorio}

La configuraci\'on del repositorio local se realiz\'o inicializando el control de
versiones en el directorio del proyecto, agregando los archivos al \'area de
preparaci\'on y registrando el commit inicial con los siguientes comandos:

\begin{lstlisting}
git init
git add .
git commit -m "Primer commit: proyecto base"
\end{lstlisting}

Posteriormente se vincul\'o el repositorio local con el repositorio remoto de
GitHub y se sincronizaron los cambios con el origen remoto mediante los comandos:

\begin{lstlisting}
git remote add origin https://github.com/Joaquinfr87/proyecto-movil.git
git push -u origin main
\end{lstlisting}

El repositorio incluye adem\'as el archivo \texttt{.gitignore} generado por
Expo, que excluye del control de versiones los directorios \texttt{node\_modules}
y otros artefactos temporales de compilaci\'on. En la figura se presenta la
evidencia de la configuraci\'on del repositorio en la terminal.

\begin{figure}[H]
  \centering
  \caption{Configuraci\'on del repositorio Git y GitHub en la terminal}
  \label{fig:configuracion-git}
  \vspace{6pt}
  \begin{lstlisting}[style=terminal]
$ git init
Initialized empty Git repository in lugares-interactivos/.git/
$ git add .
$ git commit -m "Primer commit: proyecto base"
[main 9d5c7a1] Primer commit: proyecto base
 18 files changed, 320 insertions(+)
 create mode 100644 App.tsx
 create mode 100644 app.json
 create mode 100644 package.json
 create mode 100644 pnpm-lock.yaml
 create mode 100644 tsconfig.json
$ git remote add origin https://github.com/Joaquinfr87/proyecto-movil.git
$ git push -u origin main
Enumerating objects: 20, done.
Writing objects: 100% (20/20), 1.30 MiB | 4.20 MiB/s
To https://github.com/Joaquinfr87/proyecto-movil.git
 * [new branch]      main -> main
  \end{lstlisting}
  \par\raggedright\fuenteAPA{Elaboraci\'on propia en base a la ejecuci\'on del proyecto.}
\end{figure}

En la figura se presenta la visualizaci\'on del repositorio del proyecto alojado
en GitHub.

\begin{figure}[H]
  \centering
  \caption{Repositorio del proyecto en GitHub}
  \label{fig:repositorio-github}
  \contenidoConFuente{%
    \includegraphics[width=0.62\linewidth]{repositorio.jpeg}%
  }{Elaboraci\'on propia (captura del repositorio del proyecto).}
\end{figure}
